\section{研究心法：每天否认昨天的自己}

罗福莉说，过去半年最大的感受是“每天都在否认昨天的自己”。这不是情绪化表达，而是快速范式迁移中的研究心法：当 OpenClaw、Claude Code、Agent 后训练和多模型编排快速出现时，旧判断会以周甚至天为单位失效。团队必须具备快速更新假设、快速重做实验、快速调整资源配置的能力。

\subsection{从量化到大模型的 reward 变化}

访谈中提到，量化投资里价格就是 reward，相对清晰；而大模型赛道的 reward 更不清晰、更变化。Agent 更是如此：任务成功、用户价值、安全边界、成本效率和长期学习都可能成为 reward 的一部分。这要求研究者不只优化一个数字，而要持续定义“什么才是好”。

\begin{knowledgebox}{Reward 不清晰时怎么做研究}
当 reward 不清晰，研究不能停止，而要先构造可验证代理指标：任务完成率、人工偏好、成本、失败类型、可恢复性、安全事件、用户留存等。然后不断检查这些代理指标是否真的服务最终价值。
\end{knowledgebox}

\subsection{环境比经验更重要}

“环境比经验更重要”是本期结尾的战略判断。这里的环境包括算力、基础设施、芯片、操作系统、流量、社交、组织文化和战略资源。经验固然有用，但如果环境不给反馈、不允许试错、不提供算力、不支持跨团队协作，经验会很快变成旧范式的包袱。

\begin{importantbox}{环境是研究加速度}
经验决定初始位置，环境决定加速度。Agent 时代变化太快，谁能更快获得真实反馈、更多 rollout、更好 infra、更开放组织，谁就更可能在新范式里领先。
\end{importantbox}

\subsection{本章小结}

每天否认昨天的自己，不是没有判断，而是建立能更新判断的系统。Agent 时代的研究心法，是把价值观、实验反馈、资源环境和组织调整连成一个持续学习闭环。

